% =====================================================================
%  Quantum part - Exercise 2 : dual-click MIM versus MM
% =====================================================================

\subsection{Exercise 2 -- Dual-click MIM versus MM}
\label{sec:q2}

\paragraph{Notation and common quantities.}
The one-way propagation time over the whole link is
\begin{equation}
  t_\mathrm{link} = \frac{nL}{c} = \frac{1.5\times\SI{100e3}{\metre}}{\SI{3e8}{\metre\per\second}} = \SI{500}{\micro\second}.
\end{equation}
Following the round-based controller of the lectures, a round reserves a train of complete clock slots
(launch $i$ starts at $(i-1)\,t_\mathrm{clock}$), waits for the last classical outcome, and only then
uses the confirmed pairs and resets the memories. The feedback delays are
\begin{itemize}
  \item MIM: the photon travels $L/2$ to the midpoint BSA and the result travels $L/2$ back, so the delay is $t_\mathrm{link}$;
  \item MM: the sender photon travels $L$ to the BSA at the receiver and the result travels $L$ back to the sender, so the delay is $2t_\mathrm{link}$.
\end{itemize}
Each trial succeeds independently with probability $p = p_\mathrm{dual} = 0.005$.

% ---------------------------------------------------------------------
\paragraph{1. Round duration.}
In MIM each trial reserves one memory at each endpoint, so $N$ memories per endpoint give $N$ trials per round.
In MM each trial uses a fresh sender memory, so the number of trials equals the sender-bank size $N_\mathrm{sen}$.
\begin{align}
  \text{A (MIM, } N=2\text{):}\quad
  t_\mathrm{round,A} &= t_\mathrm{link} + N\,t_\mathrm{clock} = 500 + 2\times1 = \SI{502}{\micro\second},\\
  \text{B (MM, } N_\mathrm{sen}=2\text{):}\quad
  t_\mathrm{round,B} &= 2t_\mathrm{link} + N_\mathrm{sen}\,t_\mathrm{clock} = 1000 + 2\times1 = \SI{1002}{\micro\second},\\
  \text{C (MM, } N_\mathrm{sen}=3\text{):}\quad
  t_\mathrm{round,C} &= 2t_\mathrm{link} + N_\mathrm{sen}\,t_\mathrm{clock} = 1000 + 3\times1 = \SI{1003}{\micro\second}.
\end{align}
The clock train is negligible compared with the propagation delays, so the MM rounds last about twice as long
as the MIM round.

% ---------------------------------------------------------------------
\paragraph{2. Designs A and B.}
In both designs every successful trial has its own pair of reserved memories, so no success is ever lost
(no receiver-capacity clipping). The retained yield $Y$ is the number of successes in two independent trials:
\begin{equation}
  E[Y_\mathrm{A}] = E[Y_{2,2}] = 2p = 2\times0.005 = 0.01 \ \text{pairs per round},
\end{equation}
where $Y_{N_\mathrm{sen},N_\mathrm{rec}}$ denotes the retained yield of an MM design with $N_\mathrm{sen}$ sender
and $N_\mathrm{rec}$ receiver memories (Design B is $2+2$, Design C is $3+1$).
The average rate is the expected yield divided by the round duration:
\begin{equation}
  R_\mathrm{A} = \frac{2p}{t_\mathrm{link}+2t_\mathrm{clock}} = \frac{0.01}{\SI{502}{\micro\second}} = \SI{19.9}{\per\second},
  \qquad
  R_\mathrm{B} = \frac{2p}{2t_\mathrm{link}+2t_\mathrm{clock}} = \frac{0.01}{\SI{1002}{\micro\second}} = \SI{9.98}{\per\second}.
\end{equation}

% ---------------------------------------------------------------------
\paragraph{3. Design C.}
The three sender memories launch three trials in successive slots. The single receiver memory resets locally
after each failure and can try again, but after the first success it is preserved and cannot load another pair
in the same round. Counting the $k$ potential successes among the 3 trials (binomial distribution) and the
pairs actually retained, $\min(k,1)$:
\begin{center}
\begin{tabular}{cccc}
  \toprule
  Successes $k$ & Probability $\binom{3}{k}p^k(1-p)^{3-k}$ & Retained (1 B slot) & Without clipping \\
  \midrule
  0 & $(1-p)^3 = 0.985075$     & 0 & 0 \\
  1 & $3p(1-p)^2 = 0.014888$   & 1 & 1 \\
  2 & $3p^2(1-p) = \num{7.46e-5}$ & 1 & 2 \\
  3 & $p^3 = \num{1.25e-7}$    & 1 & 3 \\
  \bottomrule
\end{tabular}
\end{center}
With $X\sim\mathrm{Binomial}(N_\mathrm{sen},p)$ the number of potential successes and $N_\mathrm{rec}=1$,
the retained yield is $\min(X,N_\mathrm{rec})$: at least one pair is retained unless all three trials fail,
\begin{equation}
  E[Y_{3,1}] = E[\min(X,1)] = 1-(1-p)^3 = 3p - 3p^2 + p^3 = 1 - 0.995^3 = 0.014925 \ \text{pairs per round},
\end{equation}
only slightly below the unclipped value $3p = 0.015$, because two or more successes in the same round are very
unlikely at $p = 0.005$. The corresponding rate is
\begin{equation}
  R_\mathrm{C} = \frac{1-(1-p)^3}{2t_\mathrm{link}+3t_\mathrm{clock}} = \frac{0.014925}{\SI{1003}{\micro\second}} = \SI{14.9}{\per\second}.
\end{equation}

% ---------------------------------------------------------------------
\paragraph{4. Comparison.}
\begin{center}
\begin{tabular}{lcccc}
  \toprule
  Design & Memories (sender + receiver) & $t_\mathrm{round}$ & $E[Y]$ per round & Rate \\
  \midrule
  A -- MIM            & 2 + 2 (per endpoint) & \SI{502}{\micro\second}  & 0.0100   & \textbf{\SI{19.9}{\per\second}} \\
  B -- symmetric MM   & 2 + 2                & \SI{1002}{\micro\second} & 0.0100   & \SI{9.98}{\per\second} \\
  C -- asymmetric MM  & 3 + 1                & \SI{1003}{\micro\second} & 0.014925 & \SI{14.9}{\per\second} \\
  \bottomrule
\end{tabular}
\end{center}
\emph{Why moving one memory to the sender improves MM.} In MM a sender memory stays reserved for the whole
$2t_\mathrm{link}$ round trip, so the number of trials per round is set by the sender bank. The receiver,
instead, learns the BSA outcome locally and promptly, so after a failure it can reset and reuse the same memory
for the next incoming photon. With $p = 0.005$ a second receiver memory is almost always idle (it is needed only
when two trials succeed in the same round, probability $\approx p^2$). Moving it to the sender buys one extra
trial per round at the cost of only one extra clock slot (\SI{1}{\micro\second}), and the rate grows by
\begin{equation}
  \frac{R_\mathrm{C}}{R_\mathrm{B}} = \frac{1-(1-p)^3}{2p}\cdot\frac{2t_\mathrm{link}+2t_\mathrm{clock}}{2t_\mathrm{link}+3t_\mathrm{clock}}
  = 1.4925\times0.999 = 1.49 \quad(+49\%).
\end{equation}

\emph{Preferred design.} Under the four-memory constraint, Design A (MIM) gives the highest rate,
\SI{19.9}{\per\second}. With the BSA at the midpoint its feedback delay is $t_\mathrm{link}$ instead of
$2t_\mathrm{link}$, so it completes about twice as many rounds per second as either MM design, and this
advantage outweighs the extra trial of Design C. Among the MM designs, C is preferred over B.
